% Lösungen Einheit 1 von 4 – Parameterform einer Ebene (zusammenbau v0.1)
\einheitenkopf{Einheit 1 von 4 · Parameterform einer Ebene}

\erg{10}{a) drei Punkte \quad b) $E\colon \vec x = (2 | 1 | 3) + r \cdot (1 | 0 | 4) + s \cdot (0 | 3 | 1)$, $r, s \in \mathbb{R}$ \quad c) $E\colon \vec x = (1 | 0 | 3) + r \cdot (2 | 2 | 1) + s \cdot (2 | -1 | -2)$, $r, s \in \mathbb{R}$ – mit $\overrightarrow{AB}$ und $\overrightarrow{AD}$, nicht mit $\overrightarrow{AB}$ und $\overrightarrow{DC}$ (parallel) \quad d) $E\colon \vec x = (1 | 2 | 0) + r \cdot (1 | 1 | 2) + s \cdot (2 | -1 | 4)$, $r, s \in \mathbb{R}$ \quad e) $L\colon \vec x = (6 | 6 | 0) + r \cdot (0 | -3 | 3) + s \cdot (-3 | 0 | 3)$, $r, s \in \mathbb{R}$; aus I $s = 1$, aus II $r = 1$, III: $3 + 3 = 6$ ✓, $S$ liegt in $L$ \quad f) ja: gleicher Stützpunkt, Spannvektoren sind Vielfache: $(2 | 0 | 2) = 2 \cdot (1 | 0 | 1)$, $(0 | -3 | -3) = -3 \cdot (0 | 1 | 1)$ \quad g) $O$, $A$ und $C$ liegen auf einer Geraden, weil $\overrightarrow{OC}$ ein Vielfaches von $\overrightarrow{OA}$ ist; $\overrightarrow{OB}$ ist kein Vielfaches von $\overrightarrow{OA}$, also liegt $B$ nicht auf dieser Geraden. Gerade und Punkt legen genau eine Ebene fest, und sie enthält alle vier Punkte.}
\erg{11}{a) $\overrightarrow{AB}$ und $\overrightarrow{DC}$ sind parallel und spannen keine Ebene auf; richtig: $\vec x = (1 | 2 | 1) + r \cdot (6 | 0 | 0) + s \cdot (0 | 3 | 4)$ mit $\overrightarrow{AD}$, $r, s \in \mathbb{R}$ \quad b) Jeder ihrer unendlich vielen Punkte kann Stützpunkt sein, und Vielfache und Summen der Spannvektoren spannen dieselbe Ebene auf. \quad c) ja: mit $\vec x = (0 | 0 | 3) + r \cdot (2 | 0 | 0) + s \cdot (0 | 1{,}5 | 2)$ folgt $r = 0{,}5$, $s = 0{,}5$, III: $3 + 1 = 4$ ✓; der Sensor sitzt in der Mitte der Platte \quad d) z. B. $(1 | 2 | 0)$, $(4 | 2 | 1)$ und $(1 | 3 | 2)$ (mit $r$, $s$ gleich 0 oder 1)}
